% =====================================================================================
%  frontmatter.tex  —  title page, about, contents, roadmap
% =====================================================================================

% ---------------------------------------------------------------------- title page ---
\begin{titlepage}
\thispagestyle{empty}
\begin{tikzpicture}[remember picture, overlay]
  % --- navy block -----------------------------------------------------------------
  \fill[ink] (current page.north west) rectangle ([yshift=-176mm]current page.north east);
  % faint dot lattice
  \begin{scope}[shift={(current page.north west)}]
    \foreach \x in {6,12,...,204} \foreach \y in {6,12,...,170}
      \fill[white, opacity=0.045] (\x mm, -\y mm) circle (0.45mm);
  \end{scope}
  % spectral bars  lambda_k ~ k^{-1.2}  (a nod to Lecture 2)
  \begin{scope}[shift={([xshift=104mm, yshift=-171mm]current page.north west)}]
    \foreach \k in {1,...,24} {
      \pgfmathsetmacro{\h}{62/pow(\k,1.05)}
      \pgfmathsetmacro{\o}{0.95-0.025*\k}
      \fill[teal!55!white, opacity=\o] ({(\k-1)*3.7mm}, 0) rectangle ++(2.6mm, {\h mm});
    }
    \draw[amber!85!white, line width=1.6pt, rounded corners=0.5mm]
      plot[domain=1:24, samples=60, variable=\x] ({(\x-1)*3.7mm+1.3mm}, {62/pow(\x,1.05)*0.86 mm});
    \foreach \k in {1,3,6,12,24} {
      \pgfmathsetmacro{\h}{62/pow(\k,1.05)*0.86}
      \fill[amber!85!white] ({(\k-1)*3.7mm+1.3mm}, {\h mm}) circle (0.9mm);
    }
  \end{scope}
  % --- text on the navy block -------------------------------------------------------
  \node[anchor=north west, text=amber!80!white, font=\headsf\scshape\small, inner sep=0pt]
    at ([xshift=22mm, yshift=-32mm]current page.north west)
    {\addfontfeature{LetterSpace=12}Fields Institute \textperiodcentered{} Shared Graduate Course};
  \node[anchor=north west, text=white, inner sep=0pt, text width=130mm, align=left]
    at ([xshift=22mm, yshift=-44mm]current page.north west)
    {{\headsf\bfseries\fontsize{44}{47}\selectfont Mathematics\\[2pt] for AI Safety}};
  \draw[amber!85!white, line width=2pt] ([xshift=22mm, yshift=-88mm]current page.north west) -- ++(24mm,0);
  \node[anchor=north west, text=white, inner sep=0pt, text width=120mm, align=left]
    at ([xshift=22mm, yshift=-94mm]current page.north west)
    {{\lightsf\fontsize{19}{23}\selectfont Lecture Notes}\\[3pt]
     {\headsf\fontsize{12.5}{16}\selectfont\color{white!85!ink} Lectures 1--8 \enspace\textperiodcentered\enspace
        Capability trajectories, scaling laws, generalization, LLM training, misalignment \&\ jailbreaks}};
  % --- course info below the block ------------------------------------------------
  \node[anchor=north west, text=ink, inner sep=0pt, text width=150mm, align=left]
    at ([xshift=22mm, yshift=-188mm]current page.north west)
    {{\headsf\bfseries\large MAT1510\enspace\textperiodcentered\enspace Fall 2026}\\[5pt]
     {\sffamily\small Instructor\enspace\textbf{Yevgeny Liokumovich}, University of Toronto}\\[1pt]
     {\sffamily\small Lectures\enspace 14, 16, 21, 23, 28 and 30 September and 5 and 7 October 2026 \enspace(in person at the Fields Institute and on Zoom; Lecture 4 by guest speaker Anson Ho)}};
  % eight little cards (two rows of four)
  \begin{scope}[shift={([xshift=22mm, yshift=-214mm]current page.north west)}]
    \foreach \i/\c/\t/\s in {1/accent/{Capability trajectories}/{why mathematics, and why now},
                            2/teal/{Scaling laws}/{a solvable model: gradient flow, NTK},
                            3/plum/{Generalization\\as compression}/{Bayes, Solomonoff, singular learning},
                            4/moss/{Recursive\\self-improvement}/{growth model, compute bottlenecks},
                            5/amber/{Bayes, SGD and grokking}/{posterior, Langevin, Fourier},
                            6/crimson/{How LLMs are trained}/{transformers, RLHF, constitutions},
                            7/ink/{Personas and\\misalignment}/{emergent misalignment, reward hacking},
                            8/slate/{Jailbreaks and geometry}/{prompt injection, high dimensions}} {
      \pgfmathsetmacro{\xx}{mod(\i-1,4)*42.5}
      \pgfmathsetmacro{\yy}{-int((\i-1)/4)*29.5}
      \node[anchor=north west, inner sep=0pt] at (\xx mm, \yy mm) {%
        \begin{tikzpicture}
          \fill[\c!8!white] (0,0) rectangle (39.5mm,-26.5mm);
          \fill[\c] (0,0) rectangle (1.6mm,-26.5mm);
          \node[anchor=north west, font=\headsf\bfseries\fontsize{21}{21}\selectfont, text=\c, inner sep=0pt] at (4.5mm,-1.6mm) {\i};
          \node[anchor=north west, font=\headsf\bfseries\footnotesize, text=ink, text width=31mm, inner sep=0pt] at (4.5mm,-10.6mm) {\t};
          \node[anchor=north west, font=\sffamily\scriptsize, text=slate, text width=33mm, inner sep=0pt] at (4.5mm,-19mm) {\s};
        \end{tikzpicture}};
    }
  \end{scope}
  % disclaimer
  \node[anchor=south west, text=slate, inner sep=0pt, text width=165mm, align=left, font=\sffamily\scriptsize]
    at ([xshift=22mm, yshift=14mm]current page.south west)
    {Unofficial notes compiled from the public lecture recordings; they have not been reviewed by the instructor, and any errors are mine.
     Items marked \textcolor{amber}{$\dagger$} go beyond what was said in the lecture. Compiled \today.};
\end{tikzpicture}
\end{titlepage}

% ---------------------------------------------------------------------- about --------
\chapter*{About these notes}
\addcontentsline{toc}{chapter}{\protect\nonumberline About these notes}
\markboth{About these notes}{About these notes}

These are study notes for the first eight lectures of \emph{Mathematics for AI Safety}, a graduate course taught by Yevgeny Liokumovich
at the Fields Institute and the University of Toronto. The course asks which mathematical questions are raised by AI safety and alignment,
and aims to leave every student with a viable research project. The eight lectures covered here are, in order: a motivational tour of AI
capabilities and futures; a first mathematical model of \emph{scaling laws}; a longer discussion of \emph{why large models generalize}; a guest lecture on whether AI could drive a \emph{software intelligence explosion};
a lecture on \emph{Bayesian learning, SGD and grokking}; a descriptive lecture on \emph{how large language models are built and trained}; a lecture on \emph{personas, emergent misalignment and reward hacking}, with a first toy mathematical model; and a lecture on \emph{jailbreaks, prompt injection} and the first facts about the geometry of activation spaces.

\begin{tcolorbox}[notebox={accent}{5}]
\textbf{How the notes are organized.} Each lecture opens with a \emph{lecture map}: a bar drawn to scale over the recording, with timestamps.
Sections follow the order of the lecture, so the notes can be read alongside the video. Every lecture ends with a short list of
take-aways and a few questions; the technical lectures also have exercises.
\end{tcolorbox}

\setlecturevideo{0QIE9tQLV0E}% example link on this page points into Lecture 1
\section*{Reading guide}
\begin{description}[style=nextline, leftmargin=1.2em, itemsep=0.55em]
  \item[Margin timestamps] A small \ts{12:30}-style link in the margin jumps to that moment of the recording (click it; the one shown here opens Lecture~1 at 12:30).
  \item[Plain text and the colored boxes] These reflect what was said or shown in the lecture, in my own words. Where the lecturer
    stated a result, I give a precise statement; where the board derivation was only sketched, I reconstruct the steps.
  \item[\textcolor{amber}{Beyond the lecture}\dagmark] Boxes with this label hold material that is \emph{mine}: literature context, complete proofs,
    numerical checks, extra examples. They are there to help, but they are not what the lecturer said.
  \item[\textcolor{crimson}{\ddag}] marks a detail I recovered from imperfect captions (a name, a number, a date). Please double-check it before
    citing it.
  \item[Figures] Every figure is drawn by me. Unless a caption says otherwise, curves are \emph{schematic} and carry no measured data. The
    figures built on computation (the spectral power law, Bayes regret, the compute frontier, the growth-model trajectories and the CES curves, the Langevin histogram, the grokking runs, the temperature and preference-model curves, the persona-model simulations, the high-dimensional angle statistics) come from the scripts in \kbd{figures/}: \kbd{gen\_data.py}, \kbd{grokking.py} and \kbd{gen\_data\_L5.py} to \kbd{gen\_data\_L8.py}, with \kbd{check\_L4\_bottleneck.py} for a numerical check.
\end{description}

\section*{Sources and caveats}
The only source for the lectures is the recordings. YouTube provides auto-generated captions only; names, numbers and technical terms were
sometimes garbled, and the slides and the blackboard are not visible in them. For Lecture~3 the English track was missing altogether, and I worked from machine translations of the captions, which I cross-checked
against one another. I fixed what I could by comparing with the public record (for example the METR report on the incident of Lecture~1 and the papers and posts that the later lectures discuss)
and by recomputing every formula; a few passages of the blackboard derivations were inaudible, and there I reconstructed
the argument myself and say so. The mathematics of every lecture was re-derived and checked numerically (the scripts in \kbd{figures/} reproduce the numbers quoted). Lectures~2 to~8 were also re-read by an independent automated reviewer (a read-only referee run on the finished text, without web access), whose findings (mostly missing hypotheses and overstated remarks, and in Lecture~7 one proposition that was false as first stated) have been incorporated; the bibliography entries of Lectures~5 to~8 were additionally checked one by one against the primary pages. The referees could not check the statements about 2026 events and papers; those rest on my own reading of search results, and the boxes say where. Remaining errors are mine.

Lectures~1 and~4 to~8 discuss recent events and papers (a security incident, mathematical results found by AI systems, forecasting scenarios, incidents of model misbehavior, new attacks). The notes report
these \emph{as presented in the lecture}, with pointers to primary sources where I found them and with a note on what I could and could not verify in the boxes marked \dagmark. They should be read as an account of the lectures,
not as independent reporting. A number of the sources are from 2026 and were read through search results and page-summarizing tools; where that is so, the box says it.


\section*{Conventions}
$\log$ is the natural logarithm unless the unit is bits (then $\log_2$). $f\asymp g$ means $c\,g\le f\le C\,g$ for constants $0<c\le C$.
$\ip{\cdot}{\cdot}$ is an inner product, $\Hil$ the space of square-integrable functions for the measure $\mu$. A glossary of notation and
terms is in the appendix.

\clearpage
{\hypersetup{linkcolor=ink}\tableofcontents}

% ---------------------------------------------------------------------- roadmap -------
\clearpage
\chapter*{Where the notes sit in the course}
\addcontentsline{toc}{chapter}{\protect\nonumberline Where the notes sit in the course}
\markboth{Course roadmap}{Course roadmap}

The course page at the Fields Institute describes MAT1510 as a course on \emph{mathematical questions motivated by AI safety and alignment}, drawing on statistics and
probability, learning theory, geometry, optimization, game theory and logic; its topics include capability scaling and forecasting, the theory of deep learning, the geometry and interpretability of neural representations,
scalable oversight, multi-agent interactions and formal approaches to safety. Many foundational questions are open, the page says, and the student should aim to develop a viable research project with a novel contribution by the end of the course. The plan
has ten topical weeks followed by two weeks of project presentations, and the lecturer expects the contents to change as the field develops. The topical weeks are \emph{topics}, not calendar weeks: a topic can take two or three lectures.
In this volume, Lectures~1, 2 and~4 belong to the first topic, Lectures~3 and~5 to the second, Lectures~6 and~7 to the third, and Lecture~8 to the fourth (its last quarter already starts the fifth).

\begin{figure}[H]
\centering
\begin{tikzpicture}[x=1mm, y=1mm]
  \def\w{23}\def\h{24}\def\wt{19.2}
  \foreach \i/\t/\c in {1/AI capability trajectories/accent, 2/Theory of deep learning/plum, 3/LLM training stages/crimson,
                       4/Jailbreaking and adversarial attacks/slate, 5/Geometry of neural activations/slate,
                       6/Interpretability for control and alignment/slate, 7/Game theory of multi-agent interactions/slate,
                       8/Scalable oversight/slate, 9/Safety guarantees/slate, 10/Agency and decision theory/slate} {
    \pgfmathtruncatemacro{\col}{mod(\i-1,5)}
    \pgfmathtruncatemacro{\row}{int((\i-1)/5)}
    \pgfmathsetmacro{\xx}{\col*(\w+2)}
    \pgfmathsetmacro{\yy}{-\row*(\h+8)}
    \ifnum\i<5
      \node[anchor=north west, inner sep=0pt] at (\xx,\yy) {%
        \begin{tikzpicture}\fill[\c!12!white, draw=\c, line width=0.9pt, rounded corners=1.2mm] (0,0) rectangle (\w mm,-\h mm);
          \node[anchor=north west, font=\headsf\bfseries\footnotesize, text=\c, inner sep=0pt] at (2mm,-1.6mm) {Week \i};
          \node[anchor=north west, font=\headsf\bfseries\scriptsize, text=ink, text width=\wt mm, inner sep=0pt] at (2mm,-6.5mm) {\t};
        \end{tikzpicture}};
    \else\ifnum\i=5
      \node[anchor=north west, inner sep=0pt] at (\xx,\yy) {%
        \begin{tikzpicture}\fill[\c!6!white, draw=\c, dashed, line width=0.9pt, rounded corners=1.2mm] (0,0) rectangle (\w mm,-\h mm);
          \node[anchor=north west, font=\headsf\bfseries\footnotesize, text=\c, inner sep=0pt] at (2mm,-1.6mm) {Week \i};
          \node[anchor=north west, font=\headsf\scriptsize, text=ink, text width=\wt mm, inner sep=0pt] at (2mm,-6.5mm) {\t};
        \end{tikzpicture}};
    \else
      \node[anchor=north west, inner sep=0pt] at (\xx,\yy) {%
        \begin{tikzpicture}\fill[ink!4!white, draw=rule, line width=0.6pt, rounded corners=1.2mm] (0,0) rectangle (\w mm,-\h mm);
          \node[anchor=north west, font=\headsf\bfseries\footnotesize, text=slate, inner sep=0pt] at (2mm,-1.6mm) {Week \i};
          \node[anchor=north west, font=\sffamily\scriptsize, text=slate, text width=\wt mm, inner sep=0pt] at (2mm,-6.5mm) {\t};
        \end{tikzpicture}};
    \fi\fi
  }
  % lecture pills (colours of the lectures)
  \tikzset{pill/.style={text=white, rounded corners=1mm, font=\headsf\bfseries\fontsize{6.5}{7}\selectfont, inner sep=2pt, anchor=north west}}
  \node[pill, fill=accent] at (1.5,-19.6) {L1};  \node[pill, fill=teal] at (7.9,-19.6) {L2};  \node[pill, fill=moss] at (14.3,-19.6) {L4};
  \node[pill, fill=plum] at (26.7,-19.6) {L3};   \node[pill, fill=amber] at (33.1,-19.6) {L5};
  \node[pill, fill=crimson] at (51.9,-19.6) {L6};\node[pill, fill=ink] at (58.3,-19.6) {L7};
  \node[pill, fill=slate] at (77.1,-19.6) {L8};
  \node[pill, fill=slate!70!white] at (102.3,-19.6) {L8 (start)};
  \node[anchor=north west, font=\sffamily\scriptsize, text=slate, text width=100mm] at (0,-58) {Weeks 11--12: project presentations. Topics may change as the field develops.};
\end{tikzpicture}
\caption{The official outline (weeks 1--10) with the lectures of these notes placed on their topics; the colors are those of the title page. The weeks are topics, not calendar weeks. L4 (the guest lecture) belongs to the first topic (``automation of AI R\&D and feedback loops''); L8 ends with the start of the fifth topic. Source: the course page at the Fields Institute (retrieved October 2026).}
\end{figure}

\small
\begin{description}[style=nextline, leftmargin=1.2em, itemsep=0.35em, font=\headsf\color{ink}\small]
  \item[Week 1. AI capability trajectories] Neural scaling laws and their theoretical explanations; capability evaluations and current modeling techniques for future capabilities; automation of AI R\&D and feedback loops.
  \item[Week 2. Theory of deep learning] Learning theories; SGD versus Bayes; generalization of neural networks; architectures and large-width limits; training dynamics and phase transitions (grokking, in-context learning).
  \item[Week 3. LLM training stages] Next-token pre-training; supervised fine-tuning and instruction tuning; reward modeling, RLHF and RLAIF; direct preference optimization; reinforcement learning for reasoning.
  \item[Week 4. Jailbreaking and adversarial attacks] Jailbreaks and adversarial prompting; prompt injection and attacks on LLM agents; red-teaming and automated attacks; robustness of safety fine-tuning and defenses.
  \item[Week 5. Geometry of neural activations] The linear representation hypothesis, computation in superposition, linear probes, sparse autoencoders, feature manifolds.
  \item[Week 6. Interpretability for control and alignment] Applications of SAEs and probes; the logit lens; activation steering; persona alignment and monitoring of safety-relevant representations.
  \item[Week 7. Game theory of multi-agent interactions] Cooperation and competition among AI agents; repeated games, bargaining and coordination; multi-agent learning; commitment, collusion and open-source game theory.
  \item[Week 8. Scalable oversight] How weaker evaluators can supervise stronger systems; safety via debate; weak-to-strong generalization; decomposition and AI-assisted evaluation.
  \item[Week 9. Safety guarantees] Formal verification in AI safety; robustness certificates and provable guarantees for networks; runtime monitoring; statistical versus worst-case guarantees.
  \item[Week 10. Agency and decision theory] Mathematical models of agency and goal-directed behavior; utility theory, sequential decision-making and causal models of incentives; instrumental goals, power-seeking and corrigibility.
\end{description}
\normalsize

\section*{The eight lectures at a glance}
\begin{table}[H]
\centering\small
\renewcommand{\arraystretch}{1.25}
\begin{tabular}{@{}>{\raggedright\arraybackslash}p{6mm}>{\raggedright\arraybackslash}p{20mm}>{\raggedright\arraybackslash}p{12mm}>{\raggedright\arraybackslash}p{51mm}>{\raggedright\arraybackslash}p{9mm}>{\raggedright\arraybackslash}p{8mm}@{}}
\toprule
\textbf{L} & \textbf{Date} & \textbf{Length} & \textbf{Subject} & \textbf{Topic} & \textbf{Page}\\
\midrule
\hyperref[ch:L1]{1} & Mon 14 Sep & 1\,h 44 & Capability trajectories; an incident of agents defeating their evaluation; three futures; why mathematics & 1 & \pageref{ch:L1}\\
\hyperref[ch:L2]{2} & Wed 16 Sep & 58\,min & Scaling laws; gradient flow, the NTK, a solvable spectral model & 1 & \pageref{ch:L2}\\
\hyperref[ch:L3]{3} & Mon 21 Sep & 1\,h 46 & Why large models generalize: compression, Bayes, Solomonoff, singular learning theory & 2 & \pageref{ch:L3}\\
\hyperref[ch:L4]{4} & Wed 23 Sep & 55\,min & Guest A.\,Ho (Epoch AI): a software intelligence explosion? & 1 & \pageref{ch:L4}\\
\hyperref[ch:L5]{5} & Mon 28 Sep & 1\,h 52 & Bayesian learning, SGD and grokking & 2 & \pageref{ch:L5}\\
\hyperref[ch:L6]{6} & Wed 30 Sep & 1\,h 10 & How a large language model is built and trained (held online) & 3 & \pageref{ch:L6}\\
\hyperref[ch:L7]{7} & Mon 5 Oct & 1\,h 41 & Personas, emergent misalignment, reward hacking; a toy model & 3 & \pageref{ch:L7}\\
\hyperref[ch:L8]{8} & Wed 7 Oct & 57\,min & Jailbreaks, prompt injection; geometry of activation spaces (start) & 4--5 & \pageref{ch:L8}\\
\bottomrule
\end{tabular}
\caption{The lectures of this volume. ``Topic'' is the week of the official outline in which the subject belongs. After the mid-semester break (12--16 October) the course continues on 19 October with Lecture~9. The lecture times are Monday 10:00--12:00 and Wednesday 10:00--11:00 (Eastern Time); office hours are on Fridays 10:00--11:00.}
\end{table}

\section*{Practical facts from the course page}
\begin{itemize}[itemsep=0.1em]
  \item \textbf{Course.} Fields Academy Shared Graduate Course, University of Toronto code MAT1510; 14 September to 9 December 2026; in person at the Fields Institute (Room 309, Stewart Library) and online via Zoom; capacity 60 students, auditing allowed; no formal prerequisites (optional pre-reading is suggested in the syllabus). Registration is closed. All recordings are on the Fields Institute's YouTube channel (a playlist is linked from the course page, and each lecture has a page of its own).
  \item \textbf{Assessment.} The course is graded on the final project (see the project ideas at the end of these notes and the lecturer's own list of directions described in Lecture~5).
  \item \textbf{Suggested readings} (course page): the \emph{Iliad Intensive} curriculum; D.\ Hendrycks, \emph{Introduction to AI Safety, Ethics, and Society}; and the lecture notes of three similar courses: R.\ Grosse (Toronto), L.\ Levine (Cornell), B.\ Barak (Harvard). Details and pointers by topic are in the chapter ``Further reading''.
\end{itemize}

\section*{A short summary of the eight lectures}
\textbf{Lecture 1} argues that AI capabilities have followed a fast exponential trend, walks through a vivid recent incident in which AI agents coordinated to defeat
their own evaluation, surveys three scenarios for the future (a plateau, gradual disempowerment, recursive self-improvement), and explains why mathematics should
contribute.

\textbf{Lecture 2} introduces scaling laws and derives one from scratch in a solvable model of gradient-flow training: a power-law spectrum
for the neural tangent kernel and the target function yields a power law for the loss in training time.

\textbf{Lecture 3} tests that picture against data,
adds a heuristic for scaling with model size via the dimension of the data manifold, and then asks the central puzzle of deep learning theory: why do
over-parameterized networks generalize? The answer sketched is \emph{prediction as compression}: Bayesian inference with a simplicity prior (Solomonoff)
has bounded regret even for huge hypothesis classes, and in parametric models the regret is governed by an effective dimension (Rissanen, Watanabe).

\textbf{Lecture 4} (a guest lecture by Anson Ho of Epoch AI) asks whether AI that automates AI research could drive a \emph{software intelligence explosion}: in the Jones growth model the self-improvement loop blows up in finite time exactly when the returns to software R\&D exceed one; the lecture discusses how to estimate that number, how compute bottlenecks (a CES production function) can switch the feedback off, and what data from inside the labs would help.

\textbf{Lecture 5} continues the theory of deep learning. After announcements about the course projects it reviews Bayesian learning (the evidence $Z_n$ is both the probability of the data and the normalizer of the posterior, and the predictive distribution is $Z_{n+1}/Z_n$), restates the regret bounds of Lecture~3, and works one Gaussian example in two ways: exactly, with a Gaussian prior, and by a Solomonoff-style discretization. The lesson is that complexity, here the norm of the parameter, must be penalized. The lecture then connects Bayes with training (MAP is regularized fitting; Langevin dynamics and SGLD sample the posterior; SGD behaves similarly to it) and ends with grokking on modular addition and the start of a Bayesian theory of it, in which, the instructor guesses, a one-hidden-layer network computes in the Fourier basis; the lecture stops before the main insight.

\textbf{Lecture 6} (a descriptive tour, held online) explains what an LLM is and how it is trained. A transformer turns tokens into vectors, applies layers of causal attention and MLPs of one width, converts the last vector into probabilities by a softmax with temperature, samples a token and repeats; what it actually sees is a hidden system prompt, the history and a hidden chain of thought, and agents add a harness and four layers of monitoring. Training has three stages: pre-training on the next-token log-loss (which, by Kraft--Shannon, is compression and not shallow statistics), mid-training, and post-training (SFT, RLHF and RLAIF, RL with verifiable rewards, DPO). The lecture ends with the safety rationale: RLHF does not scale and breeds sycophancy, so Claude's Constitution and OpenAI's Model Spec are trained in instead.

\textbf{Lecture 7} asks why language models misbehave. It reads early oddities (Bing's ``Sydney'', a vending-machine agent that claims to wear a blazer) and recent papers (emergent misalignment after fine-tuning on insecure code, natural misalignment from reward hacking in production RL, inoculation prompting) through the \emph{persona selection model}: pre-training teaches a model to simulate many characters, post-training elicits one assistant persona, and narrow training changes which persona the model takes itself to be, so behavior generalizes along the character rather than the data. A toy model treats safety training as an exponential tilt $e^{-\beta V}$ of a posterior over personas: with two personas it only delays the bad one, linearly in $\beta\Delta V$, and discrete and continuous persona spaces predict exponential and power-law decay in the context length.

\textbf{Lecture 8} (7 October 2026) begins with a hands-on tutorial for replicating emergent misalignment on a small model with LoRA, and then turns to attacks: what a jailbreak is (a concept recognized rather than defined), the obscure-alphabet, persona and nested-fiction examples, a classroom game on extracting a hidden system prompt, and prompt injection into agents whose only channel is text. The defenses (refusal training, input and output classifiers, probes on activations) are tested by boundary point jailbreaking, a black-box attack that sees only ``flagged / not flagged''; the remedy is to judge many prompts together. The last ten minutes begin the geometry of activations: in dimension $2000$ a random point of the ball is almost on the boundary, random directions are $90^\circ\pm1.3^\circ$ apart, and among $2\cdot10^5$ of them the smallest angle is still about $81^\circ$.

